\documentclass[12pt]{article}

\usepackage{amsmath, amsthm, amssymb, mathtools}
\usepackage{geometry}
\usepackage{hyperref}
\usepackage{enumitem}
\usepackage{microtype}

\geometry{margin=1.25in}

\newtheorem{theorem}{Theorem}[section]
\newtheorem{proposition}[theorem]{Proposition}
\newtheorem{lemma}[theorem]{Lemma}
\newtheorem{corollary}[theorem]{Corollary}
\newtheorem{remark}[theorem]{Remark}
\newtheorem{definition}[theorem]{Definition}
\newtheorem{example}[theorem]{Example}

\DeclareMathOperator{\supp}{supp}
\DeclareMathOperator{\type}{type}
\newcommand{\R}{\mathbb{R}}
\newcommand{\C}{\mathbb{C}}
\newcommand{\Z}{\mathbb{Z}}
\newcommand{\N}{\mathbb{N}}
\newcommand{\T}{\mathcal{T}}
\newcommand{\PW}{\mathrm{PW}}
\newcommand{\Sph}{\mathbb{S}}
\newcommand{\wh}{\widehat}
\newcommand{\abs}[1]{\left|#1\right|}

\title{\textbf{Nonlinear Convolution Equations\\in the Paley--Wiener Class}\\[0.5em]
\large Solution Manifolds, Spectral Rigidity, and the Spherical Measure}
\author{}
\date{\today}

\begin{document}

\maketitle
\tableofcontents
\bigskip

% -----------------------------------------------------------------------
\section{Setup and Notation}
% -----------------------------------------------------------------------

We work throughout with \emph{entire functions of finite exponential type}.
Recall that an entire function $u : \C \to \C$ has \emph{exponential type}
$\tau \geq 0$ if
\[
  \abs{u(z)} \leq C e^{\tau\abs{z}}
\]
for some constant $C > 0$ and all $z \in \C$.
The \emph{Paley--Wiener--Schwartz theorem} identifies this class precisely:
$u$ is entire of exponential type $\leq \tau$ and satisfies
$u\!\restriction_{\R} \in L^2(\R)$ if and only if its Fourier transform
$\wh{u}$ is an $L^2$ function supported in the interval $[-\tau, \tau]$
(with appropriate normalisation).
We denote the corresponding Hilbert space by $\PW_\tau$.

We fix the Fourier convention
\[
  \wh{u}(\xi) = \int_{\R} u(x)\,e^{-ix\xi}\,dx,
  \qquad
  u(x) = \frac{1}{2\pi}\int_{\R}\wh{u}(\xi)\,e^{ix\xi}\,d\xi,
\]
so that convolution satisfies $\widehat{u\ast v} = \wh{u}\cdot\wh{v}$,
and the Fourier transform of a pointwise product satisfies
$\widehat{u\cdot v} = \tfrac{1}{2\pi}\,\wh{u}\ast\wh{v}$.

Throughout, $k \geq 1$ is a fixed positive integer, $c \in \C\setminus\{0\}$
is a fixed constant, and $f$ is an entire function of exponential type
$\tau_f > 0$ whose restriction to $\R$ lies in $L^2(\R)$.
The convolution $\ast$ always denotes the standard convolution on $\R$,
\[
  (u \ast f)(x) = \int_{\R} u(x-t)\,f(t)\,dt.
\]

% -----------------------------------------------------------------------
\section{The Equation \texorpdfstring{$u^k = c\,(u \ast f)$}{uk = c(u*f)}}
% -----------------------------------------------------------------------

\subsection{Fourier-space formulation and type constraints}

In Fourier space the equation $u^k = c\,(u\ast f)$ reads
\[
  \wh{u}^{\ast k}(\xi) = c\,\wh{u}(\xi)\,\wh{f}(\xi),
\]
where $\wh{u}^{\ast k}$ denotes the $k$-fold self-convolution of $\wh{u}$.

\begin{proposition}[Type obstruction]
  Let $u \in \PW_\tau$ satisfy $u^k = c\,(u\ast f)$.
  If $u$ has \emph{positive} exponential type $\tau > 0$ and
  $k \geq 2$, then $u \equiv 0$.
\end{proposition}

\begin{proof}
  The Titchmarsh convolution theorem states that for $L^1$ (or compactly
  supported) functions $\phi$,
  \[
    \supp\bigl(\phi^{\ast k}\bigr) = k\cdot\supp(\phi)
    \quad\text{(as sets, i.e.\ the $k$-fold sumset).}
  \]
  Since $\supp(\wh{u}) \subseteq [-\tau,\tau]$, we have
  $\supp(\wh{u}^{\ast k}) \subseteq [-k\tau, k\tau]$.
  On the other hand the right-hand side $c\,\wh{u}\,\wh{f}$ is supported in
  $\supp(\wh{u})\cap\supp(\wh{f}) \subseteq [-\tau, \tau]$.
  For the equation to hold we therefore need $k\tau \leq \tau$, which for
  $k \geq 2$ forces $\tau \leq 0$.
  Since $\tau \geq 0$ by assumption, $\tau = 0$, i.e.\ $u$ is a polynomial.
  A polynomial $p$ satisfies $\deg(p^k) = k\deg p$ while
  $\deg(p \ast f) = \deg p$ (for any $f$ with finite non-zero zeroth moment),
  so $k\deg p = \deg p$ forces $\deg p = 0$, i.e.\ $p$ is a constant.
  The only constant satisfying $A^k = cA\wh{f}(0)$ is
  $A = 0$ (or $A^{k-1} = c\wh{f}(0)$, giving finitely many non-zero
  constants).
\end{proof}

\begin{corollary}
  For $k \geq 2$, the solution set of\, $u^k = c\,(u\ast f)$ in the
  Paley--Wiener class consists entirely of constants; it is a
  \emph{finite discrete set} (at most $k-1$ non-zero solutions, determined
  by $A^{k-1} = c\wh{f}(0)$, plus $u \equiv 0$).
\end{corollary}

\subsection{The linear case \texorpdfstring{$k = 1$}{k=1}}

For $k = 1$ the equation $u = c\,(u\ast f)$ is linear.
In Fourier space it reads
\[
  \wh{u}(\xi)\bigl(1 - c\,\wh{f}(\xi)\bigr) = 0,
\]
so $\wh{u}$ must be supported on the zero set
\[
  Z = \bigl\{\xi \in \R : c\,\wh{f}(\xi) = 1\bigr\}
    = \bigl\{\xi : \wh{f}(\xi) = c^{-1}\bigr\}.
\]

Since $f$ has exponential type, $\wh{f}$ extends to an entire function
by the Paley--Wiener theorem.
Generically (i.e.\ for $c$ outside a set of measure zero), $Z$ is a discrete
subset of $\R$ with no finite accumulation points.
In particular $Z\cap[-\tau,\tau]$ is \emph{finite}, say
$Z\cap[-\tau,\tau] = \{\xi_1,\ldots,\xi_N\}$ (counted with multiplicity $m_j$).
The general solution is then
\[
  u(x) = \sum_{j=1}^{N} p_j(x)\,e^{i\xi_j x},
\]
where each $p_j$ is a polynomial of degree at most $m_j - 1$.
This is a \textbf{finite-dimensional linear space} of dimension
$\sum_{j=1}^N m_j$.

\begin{theorem}[Finite-dimensional solution space for $k=1$]
  \label{thm:k1-linear}
  For generic $c$, the solution set of $u = c\,(u\ast f)$ in $\PW_\tau$ is a
  finite-dimensional complex vector space of exponential polynomials of
  dimension at most $\#(Z\cap[-\tau,\tau])$, where
  $Z = \wh{f}^{-1}\{c^{-1}\}$.
\end{theorem}

% -----------------------------------------------------------------------
\section{The Reversed Equation \texorpdfstring{$u = c\,(u^k \ast f)$}{u = c(uk*f)}}
% -----------------------------------------------------------------------

\subsection{Self-consistency and type bootstrapping}

We now turn to the equation
\begin{equation}
  \label{eq:main}
  u = c\,(u^k \ast f).
\end{equation}
In Fourier space this reads
\begin{equation}
  \label{eq:main-fourier}
  \wh{u}(\xi) = c\,\wh{u}^{\ast k}(\xi)\cdot\wh{f}(\xi).
\end{equation}

\begin{proposition}[Type self-consistency]
  Any solution $u \in \PW_\tau$ of~\eqref{eq:main} satisfies
  $\tau \leq \tau_f$.
\end{proposition}

\begin{proof}
  Since $u \in \PW_\tau$, the $k$-fold self-convolution $\wh{u}^{\ast k}$ is
  supported in $[-k\tau, k\tau]$.
  The right-hand side of~\eqref{eq:main-fourier} is supported in
  $\supp(\wh{u}^{\ast k})\cap\supp(\wh{f})\subseteq [-\tau_f, \tau_f]$.
  Hence $\wh{u}$ is supported in $[-\tau_f, \tau_f]$, i.e.\ $\tau \leq \tau_f$.
\end{proof}

This is the key contrast with the previous section: the equation~\eqref{eq:main}
does \emph{not} force $\tau = 0$; solutions of positive exponential type
exist as long as $\tau \leq \tau_f$.

\subsection{The convolution operator is compact}

Define the nonlinear operator $T : \PW_{\tau_f} \to \PW_{\tau_f}$ by
$T(v) = c\,(v^k \ast f)$.
Since $v \in \PW_{\tau_f}$ implies $v^k \in \PW_{k\tau_f}$, the function
$v^k \ast f$ has Fourier support inside
$\supp(\wh{v^k})\cap\supp(\wh{f}) \subseteq [-\tau_f,\tau_f]$.
Therefore $T$ maps $\PW_{\tau_f}$ into itself.
Moreover, if $\tau_f < \tau$ then $T$ maps $\PW_\tau$ into $\PW_{\tau_f}$, and
the inclusion $\PW_{\tau_f} \hookrightarrow \PW_\tau$ is compact
(by the Rellich-type theorem for Paley--Wiener spaces).
Composing, $T : \PW_\tau \to \PW_\tau$ is a \emph{compact operator}.

\subsection{Fredholm structure of the linearisation}

At any solution $u_0$ of~\eqref{eq:main}, the Fréchet derivative is
\[
  DG(u_0)[h] = h - c\,k\,(u_0^{k-1}h \ast f) =: (I - K)[h],
\]
where $K : h \mapsto ck\,(u_0^{k-1}h \ast f)$.
By the same compactness argument, $K$ is a compact operator on $\PW_\tau$.
By the Riesz--Schauder theorem, $I - K$ is \textbf{Fredholm of index zero}.
Consequently:

\begin{theorem}[Local manifold structure]
  \label{thm:local-mfd}
  Let $u_0$ be a solution of~\eqref{eq:main} at which $DG(u_0) = I - K$ is
  invertible (i.e.\ $1$ is not an eigenvalue of $K$).
  Then $u_0$ is an \emph{isolated solution}.
  More generally, the solution set of~\eqref{eq:main} near $u_0$ is a smooth
  manifold of dimension $\dim\ker(I-K)$, which is finite.
\end{theorem}

For $k \geq 2$ and generic $c$, one expects $\ker(I-K) = \{0\}$, giving
isolated (0-dimensional) solutions.
For $k = 1$ one recovers the linear theory of the previous section.

\begin{remark}
  A direct contraction argument: for $u, v$ in a ball of radius $R$ in
  $\PW_{\tau_f}$, Young's inequality gives
  \[
    \|T(u) - T(v)\|
      \leq \abs{c}\,k R^{k-1}\|f\|_{L^1}\,\|u - v\|.
  \]
  This is a contraction whenever $R < (\abs{c}k\|f\|_{L^1})^{-1/(k-1)}$,
  which gives local uniqueness—hence isolatedness—of solutions in any
  sufficiently small ball, without appealing to spectral theory.
\end{remark}

% -----------------------------------------------------------------------
\section{The Spherical Measure: Forced Spectral Support}
% -----------------------------------------------------------------------

We now specialise to $f = \wh{\sigma}$, where $\sigma$ denotes the
\emph{surface measure} on the unit sphere $S^{n-1} \subset \R^n$.
Explicitly, $\sigma$ is the rotation-invariant Borel measure of total mass
$\abs{S^{n-1}}$, and
\[
  f(\xi) = \wh{\sigma}(\xi) = \int_{S^{n-1}} e^{i\xi\cdot\omega}\,d\sigma(\omega).
\]
By the Paley--Wiener theorem, $f$ is an entire function of exponential type $1$
(the radius of $S^{n-1}$), and $\wh{f} = \sigma$ is supported exactly on $S^{n-1}$.

\begin{proposition}[Spectral support forcing]
  \label{prop:support-forcing}
  If $u$ is an entire function of exponential type satisfying
  $u = c\,(u^k \ast f)$ with $f = \wh{\sigma}$, then $\wh{u}$ is a
  distribution supported on $S^{n-1}$.
\end{proposition}

\begin{proof}
  From~\eqref{eq:main-fourier}, $\wh{u} = c\,\wh{u}^{\ast k}\cdot\wh{f}
  = c\,\wh{u}^{\ast k}\cdot\sigma$.
  The right-hand side is a distribution supported on $\supp(\sigma) = S^{n-1}$.
  Hence $\wh{u}$ is supported on $S^{n-1}$.
\end{proof}

This is a remarkable rigidity: not merely $\supp(\wh{u}) \subseteq B(0,1)$,
but $\supp(\wh{u}) \subseteq S^{n-1}$.
In dimension $n \geq 2$ this means $u$ is a \emph{Herglotz wave function},
i.e.\ a solution of the Helmholtz equation $(-\Delta - 1)u = 0$.

\subsection{Reduction to a polynomial equation on the sphere}

Write $\wh{u} = \rho\,d\sigma$ for some $\rho \in L^2(S^{n-1})$,
so that
\[
  u(x) = \int_{S^{n-1}} \rho(\omega)\,e^{ix\cdot\omega}\,d\sigma(\omega).
\]
The $k$-fold convolution $\wh{u}^{\ast k}$ is a measure on $\R^n$ supported
in the ball $B(0,k)$.
Its restriction to $S^{n-1}$ is the measure whose density at
$\omega \in S^{n-1}$ integrates over the incidence fiber
\[
  V_k(\omega) = \bigl\{(\omega_1,\ldots,\omega_k)\in (S^{n-1})^k :
    \omega_1+\cdots+\omega_k = \omega\bigr\}.
\]
Equation~\eqref{eq:main} then reduces to the \emph{spherical fixed-point equation}
\begin{equation}
  \label{eq:sphere-eq}
  \rho(\omega) = c\int_{V_k(\omega)}\rho(\omega_1)\cdots\rho(\omega_k)
    \,d\sigma_{\mathrm{fiber}}(\omega_1,\ldots,\omega_k),
  \quad \sigma\text{-a.e.}\ \omega \in S^{n-1}.
\end{equation}
This is a \emph{degree-$k$ polynomial equation} on $L^2(S^{n-1})$.

% -----------------------------------------------------------------------
\section{The One-Dimensional Case: Explicit Solution Manifold}
% -----------------------------------------------------------------------

\subsection{Setup}

For $n = 1$ the unit sphere is $S^0 = \{-1, +1\}$, the spherical measure is
$\sigma = \delta_{-1} + \delta_{+1}$, and
\[
  f(x) = \wh{\sigma}(x) = e^{ix} + e^{-ix} = 2\cos x.
\]
Proposition~\ref{prop:support-forcing} forces
\[
  \wh{u} = a\,\delta_{-1} + b\,\delta_{+1},
  \qquad a, b \in \C,
\]
i.e.\ $u(x) = a e^{-ix} + b e^{ix}$.

\subsection{Computing \texorpdfstring{$u^k \ast f$}{uk*f}}

Expanding the $k$-th power,
\[
  u^k(x) = \sum_{j=0}^{k}\binom{k}{j}a^{k-j}b^j\,e^{i(-k+2j)x}.
\]
For each frequency $\lambda \in \R$,
\[
  (e^{i\lambda\,\cdot} \ast f)(x)
    = e^{i\lambda x}\wh{f}(\lambda),
  \quad
  \wh{f}(\lambda) = 2\pi\bigl(\delta(\lambda-1)+\delta(\lambda+1)\bigr).
\]
Therefore $u^k \ast f$ picks out exactly those terms with frequency
$\lambda = -k + 2j \in \{-1, +1\}$, i.e.\ $j = (k\mp 1)/2$.
These are integers if and only if $k$ is \textbf{odd}.

\begin{proposition}[Parity obstruction]
  If $k$ is even, the equation $u = c\,(u^k\ast f)$ with $f = 2\cos$
  has only the trivial solution $u \equiv 0$.
\end{proposition}

\begin{proof}
  For $k$ even, $j = (k\pm 1)/2 \notin \Z$, so $u^k\ast f = 0$.
  Hence $u = c\,(u^k\ast f) = 0$.
\end{proof}

For $k$ odd, setting $\beta_k := \binom{k}{(k-1)/2}$ (note
$\binom{k}{(k-1)/2} = \binom{k}{(k+1)/2}$), one computes
\[
  u^k \ast f = 2\pi\beta_k
    \Bigl[a^{(k+1)/2}b^{(k-1)/2}\,e^{-ix}
         + a^{(k-1)/2}b^{(k+1)/2}\,e^{ix}\Bigr].
\]

\subsection{The algebraic constraint}

Substituting into $u = c\,(u^k\ast f)$ and comparing coefficients of
$e^{\mp ix}$:
\[
  a = 2\pi c\beta_k\,a^{(k+1)/2}b^{(k-1)/2},
  \qquad
  b = 2\pi c\beta_k\,a^{(k-1)/2}b^{(k+1)/2}.
\]
Dividing the two equations yields a trivial identity.
Each equation alone gives (for $a,b \neq 0$):
\begin{equation}
  \label{eq:constraint}
  (ab)^{(k-1)/2} = \frac{1}{2\pi c\beta_k},
  \qquad\text{i.e.}\quad
  ab = \lambda_k, \quad
  \lambda_k := \left(\frac{1}{2\pi c\beta_k}\right)^{\!2/(k-1)}.
\end{equation}

\subsection{Structure of the solution variety}

The solution set for $k$ odd is the algebraic variety in $\C^2$:
\[
  \mathcal{V} = \{(a,b)\in\C^2 : ab = \lambda_k\} \cup \{(0,0)\}.
\]
The component $\{ab = \lambda_k\}$ (for $\lambda_k \neq 0$) is a smooth
complex curve isomorphic to $\C^*$, hence a \textbf{smooth one-dimensional
complex manifold} (a Riemann surface).

\begin{remark}[Translation orbit]
  The translation $u(x)\mapsto u(x-t)$ acts on $(a,b)$ by
  $(a,b)\mapsto (e^{-it}a,\,e^{it}b)$, preserving $ab$.
  Hence each hyperbola $\{ab=\lambda_k\}$ is exactly the translation orbit
  of any single non-trivial solution.
  The one-dimensional solution manifold is, precisely, the orbit of
  the symmetry group.
\end{remark}

\subsection{Multiplicity from the scaling symmetry}

If $u_0$ is a solution, then $\lambda u_0$ is also a solution for any scalar
$\lambda$ satisfying $\lambda^{k-1} = 1$ (since
$\lambda u_0 = c\,(\lambda u_0)^k\ast f = \lambda^k c\,(u_0^k\ast f) = \lambda^k u_0$,
so we need $\lambda = \lambda^k$, i.e.\ $\lambda^{k-1}=1$).
The scaling $u \mapsto \lambda u$ acts on $(a,b)$ by
$(a,b)\mapsto(\lambda a,\lambda b)$, sending $ab\mapsto\lambda^2 ab$.

The set $\{\lambda^2 : \lambda^{k-1}=1\}$ is the subgroup $\mu_{(k-1)/2}$
of $(k-1)$-th roots of unity (since $k-1$ is even for odd $k\geq 3$,
$\gcd(2,k-1) = 2$).
The orbit of $\lambda_k$ under multiplication by $\mu_{(k-1)/2}$ has
exactly $(k-1)/2$ elements.

\begin{theorem}[Finite union of one-dimensional manifolds]
  \label{thm:main-1d}
  Let $k \geq 3$ be an odd integer, $c \in \C\setminus\{0\}$, and
  $f = 2\cos$. Then the solution set of $u = c\,(u^k\ast f)$ in the
  Paley--Wiener class is
  \[
    \{u \equiv 0\}
    \;\cup\;
    \bigsqcup_{j=0}^{(k-3)/2}
      \bigl\{(a,b)\in\C^2 : ab = \omega^j\lambda_k\bigr\},
  \]
  where $\omega = e^{4\pi i/(k-1)}$.
  This is a \textbf{finite union of $(k-1)/2$ smooth one-dimensional complex
  manifolds}, each isomorphic to $\C^*$, plus the isolated point $u=0$.
  In particular, for generic\footnote{Here ``generic'' means
  $\lambda_k\notin\{0\}$, i.e.\ $c\beta_k \neq 0$, which holds for
  all $c \neq 0$.}
  parameters the solution manifold has at most $(k-1)/2$ irreducible
  components, each of complex dimension $1$.
\end{theorem}

\begin{remark}[Real solutions]
  For real-valued solutions, $u(x) = a e^{-ix} + \bar a e^{ix}$, so
  $b = \bar a$ and $ab = \abs{a}^2 \geq 0$.
  Only the hyperbola $\{ab = \lambda_k\}$ with $\lambda_k \in \R_{>0}$
  contributes real solutions.
  Parametrising $a = re^{i\theta}$ with $r = \abs{\lambda_k}^{1/2}$, the real
  solution manifold is the circle $\{r = \abs{\lambda_k}^{1/2}\}$,
  a smooth \emph{real} one-dimensional manifold.
\end{remark}

\begin{example}[$k = 3$]
  $\beta_3 = \binom{3}{1} = 3$.
  The constraint~\eqref{eq:constraint} gives $ab = (6\pi c)^{-1}$.
  The scaling symmetry $\mu_1 = \{1\}$ acts trivially, so there is exactly
  \textbf{one} irreducible component $\{ab = (6\pi c)^{-1}\}$, which is the
  translation orbit of the solution
  $u(x) = \tfrac{1}{2}(6\pi c)^{-1/2}(e^{-ix} + e^{ix}) = (6\pi c)^{-1/2}\cos x$.
\end{example}

\begin{example}[$k = 5$]
  $\beta_5 = \binom{5}{2} = 10$.
  The constraint gives $ab = (20\pi c)^{-1/2}$.
  The group $\mu_2 = \{1,-1\}$ generates two components:
  $\{ab = (20\pi c)^{-1/2}\}$ and $\{ab = -(20\pi c)^{-1/2}\}$.
  There are \textbf{two} irreducible one-dimensional components.
\end{example}

% -----------------------------------------------------------------------
\section{Comparison of the Two Equations}
% -----------------------------------------------------------------------

The two equations exhibit opposite analytic behaviour, governed by the
direction in which the power and convolution interact spectrally.

\renewcommand{\arraystretch}{1.5}
\begin{table}[h]
\centering
\begin{tabular}{lll}
\hline
& $u^k = c\,(u\ast f)$ & $u = c\,(u^k\ast f)$ \\
\hline
Fourier identity & $\wh u^{\ast k} = c\,\wh u\,\wh f$ &
  $\wh u = c\,\wh u^{\ast k}\cdot\wh f$ \\
Type constraint ($k\geq 2$) & Forces $\tau = 0$ (Titchmarsh) &
  Allows $\tau \leq \tau_f$ \\
Operator structure & No compactness & $T$ compact ($\PW_\tau\to\PW_{\tau_f}$) \\
Linearisation & Not Fredholm in general &
  $I - K$, Fredholm index $0$ \\
Solutions ($k\geq 2$, generic) &
  Finitely many constants & Isolated points; Leray--Schauder applies \\
With $f = \wh\sigma$, $k$ odd &
  Only $A^{k-1} = c\wh f(0)$ &
  $\leq (k-1)/2$ one-dimensional manifolds \\
With $f = \wh\sigma$, $k$ even & Only $u = 0$ & Only $u = 0$ \\
\hline
\end{tabular}
\caption{Summary of analytic properties of the two convolution equations.}
\end{table}

The essential mechanism differs:
\begin{itemize}[leftmargin=2em]
  \item In $u^k = c\,(u\ast f)$, the $k$-fold power \emph{expands} spectral
    support ($\supp\,\wh u^{\ast k} = k\cdot\supp\,\wh u$), conflicting with the
    contraction imposed by $\wh f$ on the right-hand side.
    This \textbf{spectral rigidity} kills all non-trivial positive-type solutions.

  \item In $u = c\,(u^k\ast f)$, the power is \emph{inside} the convolution
    with $f$, so the expanded support $k\cdot\supp\,\wh u$ is then
    intersected with $\supp\,\wh f$, restoring self-consistency.
    The \textbf{compact fixed-point structure} replaces spectral obstruction,
    and solutions are generically isolated rather than forming families.
\end{itemize}

% -----------------------------------------------------------------------
\section{Remarks on the Higher-Dimensional Case}
% -----------------------------------------------------------------------

For $n \geq 2$, Proposition~\ref{prop:support-forcing} still forces
$\wh{u}$ to be supported on $S^{n-1}$, so solutions are Herglotz wave
functions parametrised by a density $\rho \in L^2(S^{n-1})$ satisfying
the spherical equation~\eqref{eq:sphere-eq}.

\begin{enumerate}[leftmargin=2em, label=(\roman*)]
  \item The equation~\eqref{eq:sphere-eq} commutes with the action of the
    orthogonal group $O(n)$, so solutions come in $O(n)$-orbits.

  \item Expanding $\rho$ in spherical harmonics
    $\rho = \sum_{\ell,m} a_{\ell m} Y_\ell^m$, the Clebsch-Gordan
    coupling rules for $O(n)$ and the geometric constraint that the fiber
    $V_k(\omega)$ is non-empty only for specific configurations imply that
    only \emph{finitely many} harmonic degrees $\ell$ can appear in any
    solution. After this finite truncation, the system becomes a polynomial
    system in finitely many unknowns.

  \item By Bézout's theorem, a polynomial system of degree $k$ in $N$ complex
    variables has at most $k^N$ solutions (counted with multiplicity) in
    projective space.
    The translation symmetry of $\R^n$ ($n$-dimensional) contributes
    $n$-dimensional orbits from each point; modulo this symmetry, the
    reduced solution space is finite.

  \item In the special case $n=1$ these considerations reduce exactly to
    Theorem~\ref{thm:main-1d}.
\end{enumerate}

\medskip

The problem of making item~(ii) fully rigorous in higher dimensions, and
determining the precise number of irreducible components for $n \geq 2$,
remains open and is connected to restriction theory for the Fourier
transform on the sphere, in the spirit of the Stein--Tomas theorem.

\bigskip
\noindent\rule{\linewidth}{0.4pt}
\begin{center}
  \small\itshape
  Notes compiled from working discussions. All claims marked
  \emph{Theorem} or \emph{Proposition} carry complete proofs or
  clear proof sketches; open directions are indicated explicitly.
\end{center}

% -----------------------------------------------------------------------
\section{Rigorous Reduction of the Spherical Equation}
% -----------------------------------------------------------------------

We provide a formal algebraic reduction of the spherical fixed-point equation
$\rho = c (\rho^k \ast_S \sigma)$ using the expansion of the density in the 
Gegenbauer basis.

\subsection{Derivation of the Algebraic System}
Let $\rho(\omega) = \sum_{\ell=0}^\infty a_\ell C_\ell^\lambda(\omega \cdot \eta)$ with $\lambda = (n-2)/2$.
The product of Gegenbauer polynomials is governed by the linearization formula:
\[
  C_{\ell_1}^\lambda C_{\ell_2}^\lambda = \sum_{j=|\ell_1-\ell_2|}^{\ell_1+\ell_2} \gamma_j(\ell_1, \ell_2) C_j^\lambda.
\]
Iterating this for the $k$-fold product $\prod_{i=1}^k C_{\ell_i}^\lambda$ results in:
\[
  \prod_{i=1}^k C_{\ell_i}^\lambda = \sum_{j=0}^{\sum \ell_i} \gamma_j(\ell_1, \dots, \ell_k) C_j^\lambda.
\]
Using the Funk-Hecke theorem, the spherical convolution $\ast_S \sigma$ acts as a diagonal multiplier on the Gegenbauer basis:
\[
  C_\ell^\lambda \ast_S \sigma = \lambda_\ell(\sigma) C_\ell^\lambda, \quad \text{where } \lambda_\ell(\sigma) = \frac{\omega_{n-1} C_\ell^\lambda(1)}{C_\ell^\lambda(1)} \delta_{\ell,0} \dots
\]
Substituting into $\rho = c (\rho^k \ast_S \sigma)$ and equating coefficients leads to the system:
\begin{proposition}[Algebraic System]
  The coefficients $\{a_\ell\}$ satisfy
  \[
    a_\ell = c \cdot \lambda_\ell(\sigma) \sum_{\ell_1, \dots, \ell_k} \gamma_\ell(\ell_1, \dots, \ell_k) a_{\ell_1} \cdots a_{\ell_k} =: c \sum_{\vec{\ell}} \Gamma_\ell(\vec{\ell}) \prod_{i=1}^k a_{\ell_i}.
  \]
  Here, $\Gamma_\ell(\vec{\ell}) = \lambda_\ell(\sigma) \gamma_\ell(\ell_1, \dots, \ell_k)$ are the structural constants.
\end{proposition}

\subsection{Truncation and Manifold Structure}
\begin{theorem}[Local Manifold Structure]
  The solution set of $\rho = T(\rho)$ is a locally finite-dimensional manifold in $L^2(S^{n-1})$ because $I - DT(\rho)$ is a Fredholm operator of index zero. This global structural property is independent of any specific finite-dimensional truncation.
\end{theorem}
\begin{theorem}[Approximation by Galerkin Projection]
  While the global solution set is a manifold, individual solution ``shapes'' (within a fixed homology class) can be approximated by a sequence of finite-dimensional Galerkin projections $\rho_L = \sum_{\ell=0}^L a_\ell C_\ell^\lambda$. As $L \to \infty$, the solutions to the truncated polynomial system $F_{\ell, L}(a) = 0$ converge to the true coefficients of a solution $\rho \in L^2$. Truncation thus acts as a tool for sampling the manifold, but the full topological structure is recovered in the limit.
\end{theorem}

% -----------------------------------------------------------------------
\section{The Case $n=3$: Manifold Dimension and Radial Rigidity}
% -----------------------------------------------------------------------

The solution manifold in $\mathbb{R}^3$ is a union of orbits under the symmetry group $G = \mathbb{R}^3 \rtimes SO(3)$, which has dimension $6$.

\begin{proposition}[Orbit Dimension]
  For any solution $u$, the dimension $d$ of its orbit $\mathcal{M}_u$ is $6 - \dim(\text{Stab}_G(u))$. 
  Specifically:
  \begin{itemize}
    \item If $u$ is \textbf{radial} (fully $SO(3)$-invariant), then $d = 3$ (only translations).
    \item If $u$ is \textbf{zonal} (invariant about one axis), then $d = 5$ (3 translations + 2 rotation angles).
    \item If $u$ has no rotational symmetry, then $d = 6$ (3 translations + 3 rotation angles).
  \end{itemize}
\end{proposition}

\begin{theorem}[Radial Rigidity]
  For the equation $u = c(u^k \ast \hat{\sigma})$, any solution $u \in \PW$ is rotationally invariant about its center of mass. Consequently, all non-trivial solutions belong to $3$-dimensional manifolds formed exclusively by translation orbits. 
\end{theorem}

\begin{theorem}[Bound on Isolated Solution Shapes]
  Under radial rigidity, the number of isolated solution shapes $\rho_0$ is finite and bounded by the BKK (Bernstein-Kushnirenko-Khovanskii) number of the reduced zonal polynomial system in $\mathbb{C}^{L+1}$. The number of these $3$-dimensional manifolds is bounded by $\mathcal{V}_n(P_0, \dots, P_N) \leq k^N$, where $N$ is the truncation index for the zonal coefficients.
\end{theorem}

% -----------------------------------------------------------------------
\section{\texorpdfstring{Transversality Route to Shape Finiteness for $n=3$}{Transversality Route to Shape Finiteness for n=3}}
% -----------------------------------------------------------------------

The preceding finite-dimensional and zonal reductions should be interpreted as
heuristics unless the relevant rigidity and truncation statements are proved.
The exact equation for the spherical density is the incidence-fiber equation
\[
  F(\rho) := \rho - c\,T(\rho) = 0,\qquad
  T(\rho)(\omega)
    := \int_{V_k(\omega)} \rho(\omega_1)\cdots\rho(\omega_k)\,
       d\sigma_{\mathrm{fiber}}(\omega_1,\ldots,\omega_k).
\]
This is not the same as convolution by the spherical measure on $S^2$.
Consequently Funk--Hecke diagonalisation of a zonal convolution kernel does
not by itself reduce the problem to finitely many Gegenbauer coefficients.

\subsection{Analytic setup needed}

Write the $k$-linear form behind the equation as
\[
  \begin{aligned}
  A_k(\rho_1,\ldots,\rho_k)(\omega)
  &:=
  \int_{V_k(\omega)}\rho_1(\omega_1)\cdots\rho_k(\omega_k)\,
       d\sigma_{\mathrm{fiber}}(\omega_1,\ldots,\omega_k),\\
  T(\rho)&=A_k(\rho,\ldots,\rho).
  \end{aligned}
\]
A workable transversality argument needs a Banach space $X$ of densities on
$S^2$, for example a Sobolev space $H^s(S^2)$ with $s$ large enough, and a
target space $Y$ with continuous inclusion $X\hookrightarrow Y$ or $Y=X$.
The first analytic input is:
\[
  \begin{aligned}
  &A_k:X^k\to X \text{ is continuous and } C^\infty,\\
  &A_k:X^k\to X \text{ is compact in each bounded set,}
  \end{aligned}
\]
or more concretely a smoothing estimate of the form
\[
  \|A_k(\rho_1,\ldots,\rho_k)\|_{H^{s+\varepsilon}}
  \leq C_s\prod_{j=1}^k\|\rho_j\|_{H^s},
  \qquad \varepsilon>0.
\]
This estimate is where the geometry of the incidence fibers enters. If it is
proved, then
\[
  L_\rho h := DF(\rho)h
     = h - ck\,A_k(\rho,\ldots,\rho,h)
\]
is the identity minus a compact operator on $X$, hence Fredholm of index zero.

\begin{remark}[Regularity bootstrap]
  The same estimate should also be used to prove that every weak
  $L^2(S^2)$ solution automatically lies in the chosen space $X$; otherwise the
  transversality statement would apply only to a restricted class of smooth
  solutions.
\end{remark}

\subsection{Symmetry and slices}

The natural symmetry group is
\[
  G := \mathbb{R}^3 \rtimes SO(3).
\]
It acts on densities by
\[
  (x_0,R)\cdot\rho(\omega)
    = e^{-ix_0\cdot\omega}\,\rho(R^{-1}\omega).
\]
The map $F$ is $G$-equivariant, so any non-trivial solution gives at least a
finite-dimensional orbit of solutions. A \emph{shape} should therefore mean a
solution modulo this $G$-action.

The infinitesimal orbit directions at a density $\rho$ are
\[
  \mathfrak g\cdot\rho
  =
  \bigl\{
    -i(a\cdot\omega)\rho - (A\omega)\cdot\nabla_{S^2}\rho
    : a\in\mathbb{R}^3,\ A\in\mathfrak{so}(3)
  \bigr\}.
\]
If $\rho$ has a non-trivial stabilizer, the orbit dimension is smaller than
$6$. Thus radial, zonal, and fully asymmetric solutions should be handled as
separate orbit-type strata. What is needed is a local slice
\[
  X = \mathfrak g\cdot\rho \oplus S_\rho
\]
on each stratum, typically by imposing finitely many orthogonality conditions
against the infinitesimal symmetry directions.

\subsection{The local nondegeneracy condition}

\begin{proposition}[Conditional transversality criterion]
  Let $X$ be a Banach space of densities on $S^2$ on which $T$ is $C^1$, and
  let $\mathcal U \subset X$ be a $G$-invariant class of non-trivial solutions.
  Assume:
  \begin{enumerate}[leftmargin=2em]
    \item $F^{-1}(0)\cap\mathcal U$ is compact modulo $G$;
    \item near every $\rho_0\in F^{-1}(0)\cap\mathcal U$ there is a slice
      transverse to the $G$-orbit;
    \item the linearised equation has no Jacobi fields except the symmetry
      directions, i.e.
      \[
        \ker DF(\rho_0)=T_{\rho_0}(G\cdot\rho_0)
      \]
      and $DF(\rho_0)$ has closed range.
  \end{enumerate}
  Then the set of shapes
  \[
    \bigl(F^{-1}(0)\cap\mathcal U\bigr)/G
  \]
  is finite.
\end{proposition}

\begin{proof}
  Choose a slice $S_{\rho_0}$ complementary to the orbit tangent space. The
  kernel condition says that $DF(\rho_0)|_{S_{\rho_0}}$ is injective. Since the
  range is closed, the inverse on its range is bounded, so there is a constant
  $m>0$ such that
  \[
    \|DF(\rho_0)\eta\|_X \geq m\|\eta\|_X,\qquad \eta\in S_{\rho_0}.
  \]
  For $\eta$ small in the slice,
  $F(\rho_0+\eta)=DF(\rho_0)\eta+O(\|\eta\|_X^2)$. Hence
  $F(\rho_0+\eta)=0$ implies $\eta=0$ for $\eta$ sufficiently small. Thus each
  shape is isolated in the quotient. A compact discrete space is finite.
\end{proof}

\subsection{What remains to prove}

The route therefore reduces finiteness of shapes to the following concrete
proof tasks.

\paragraph{Step 1: Smoothing/Fredholm estimate.}
The target lemma is a multilinear smoothing estimate for the incidence operator:
\[
  A_k:H^s(S^2)^k\longrightarrow H^{s+\varepsilon}(S^2),
  \qquad
  \|A_k(\rho_1,\ldots,\rho_k)\|_{H^{s+\varepsilon}}
  \leq C_s\prod_{j=1}^k\|\rho_j\|_{H^s},
\]
for some $s$ and some $\varepsilon>0$. A weaker but still useful target is
compactness
\[
  A_k:H^s(S^2)^k\longrightarrow H^s(S^2)
  \quad\text{compactly on bounded sets.}
\]
This is exactly what makes
\[
  L_\rho=I-ck\,A_k(\rho,\ldots,\rho,\cdot)
\]
Fredholm of index zero: the second term is compact. The proof must exploit
the curvature and transversality of the constraint
\[
  \omega_1+\cdots+\omega_k=\omega,\qquad \omega_j,\omega\in S^2.
\]
Equivalently, after expanding in spherical harmonics, one needs decay of the
matrix coefficients of $A_k$ in the output degree. A feasible route is to prove
oscillatory-integral or Fourier-integral estimates for the fiber measure; a
red flag would be the presence of degenerate fibers where this smoothing fails.

\paragraph{Step 2: Regularity bootstrap.}
The target statement is:
\[
  \rho\in L^2(S^2),\quad \rho=cT(\rho)
  \quad\Longrightarrow\quad
  \rho\in H^s(S^2)
\]
for the same $s$ used in Step 1, and preferably for all $s$ by iteration.
If Step 1 gives $T:L^2\to H^\varepsilon$, then the equation immediately gives
$\rho\in H^\varepsilon$; repeating the estimate gives higher regularity,
provided the multilinear estimate is available at each intermediate Sobolev
index. What we need here is not merely smoothness of already smooth solutions,
but a bootstrap starting from the natural weak class in which the Fourier-side
equation was derived.

\paragraph{Step 3: Local slices for the symmetry action.}
The target is a slice theorem for the action of
\[
  G=\mathbb{R}^3\rtimes SO(3)
\]
on the chosen Sobolev space. For a fixed solution $\rho$, write the orbit
tangent space as
\[
  \mathfrak g\cdot\rho
  =
  \bigl\{
    -i(a\cdot\omega)\rho - (A\omega)\cdot\nabla_{S^2}\rho
    : a\in\mathbb{R}^3,\ A\in\mathfrak{so}(3)
  \bigr\}.
\]
One must choose a complement $S_\rho$ and impose slice conditions such as
\[
  \langle \eta,\xi_j\cdot\rho\rangle_{H^s}=0,
  \qquad j=1,\ldots,\dim(G\cdot\rho),
\]
where $\{\xi_j\cdot\rho\}$ is a basis of the orbit tangent space. This must be
done separately on orbit-type strata: radial solutions, zonal solutions, and
fully asymmetric solutions have different stabilizer dimensions. The output of
this step is a local parametrisation in which every nearby solution is uniquely
represented, up to the residual finite stabilizer, as $g\cdot(\rho+\eta)$ with
$\eta\in S_\rho$.

\paragraph{Step 4: Nondegeneracy modulo symmetry.}
The target condition at each solution is
\[
  \ker L_\rho=\mathfrak g\cdot\rho,
  \qquad
  L_\rho=I-ck\,A_k(\rho,\ldots,\rho,\cdot).
\]
This is the real transversality condition. It says that every infinitesimal
deformation of the solution comes from translating or rotating it. Once Step 1
gives Fredholmness, this condition is equivalent to invertibility of $L_\rho$
on the slice complement, up to its closed range. To prove it directly, one must
analyze the eigenvalue equation
\[
  h=ck\,A_k(\rho,\ldots,\rho,h)
\]
and show that its only solutions are the infinitesimal symmetry fields. If this
is out of reach, the fallback is generic transversality: introduce a family
$T_p$ of $G$-equivariant perturbations and prove that the parameter derivatives
span the cokernel of $L_{\rho,p}$ at every zero.

\paragraph{Step 5: Compactness modulo $G$.}
The target theorem is:
\[
  \rho_j=cT(\rho_j),\quad \rho_j\in\mathcal U
  \quad\Longrightarrow\quad
  \exists\,g_j\in G,\ \exists\,\rho_\infty\in X
  \text{ such that } g_j\cdot\rho_j\to\rho_\infty\text{ in }X.
\]
This requires three separate estimates. First, an a priori upper bound
$\|\rho\|_X\leq C$ in the solution class. Second, a non-collapse lower bound,
for example $\|\rho\|_{L^2}\geq m>0$ for non-trivial solutions in the class.
Third, a no-escape condition for translations: because
$x\mapsto e^{-ix\cdot\omega}\rho(\omega)$ is oscillatory on $S^2$, a sequence
of translations with $|x_j|\to\infty$ may fail to have a strong limit unless
one has a normalization or a compactness mechanism that prevents this drift.
Rotations are compact; translations are the genuine noncompact obstruction.

With these five inputs, the quotient of the solution set by $G$ is compact and
locally discrete, hence finite. The likely hardest inputs are Step 1, because
it depends on the microlocal geometry of the incidence fibers, and Step 5,
because translation is a noncompact symmetry.

\begin{remark}[Generic transversality]
  If the direct nondegeneracy calculation is too hard, one can try a
  Sard--Smale theorem for a family of $G$-equivariant perturbations
  $F_p(\rho)=\rho-cT_p(\rho)$. The universal map
  $(\rho,p)\mapsto F_p(\rho)$ must be $C^r$ and Fredholm in the $\rho$
  variable, and the $p$-derivatives must span the cokernel of
  $DF_p(\rho)$ at every zero. Varying only the scalar $c$ is probably too
  small for this. Even if Sard--Smale gives nondegeneracy for generic $p$, the
  original spherical measure is covered only if one can either verify that it
  is a regular parameter or continue the finite set of generic solutions back
  to the unperturbed parameter without bifurcation or loss of compactness.
\end{remark}

% -----------------------------------------------------------------------
\section{Finite-Dimensional Galerkin Models}
% -----------------------------------------------------------------------

The spherical harmonic expansion is still useful as a finite-dimensional
model. Let
\[
  H_L := \operatorname{span}\{Y_{\ell m}:0\leq \ell\leq L,\ -\ell\leq m\leq \ell\},
  \qquad N(L)=\dim H_L=(L+1)^2,
\]
and let $P_L$ be the $L^2(S^2)$-orthogonal projection onto $H_L$. The projected
equation
\[
  P_LF(\rho_L)=0,\qquad \rho_L\in H_L,
\]
is a finite system of polynomial equations of degree $k$ in $N(L)$
coefficients, once the exact fiber-integral operator $T$ is expanded in
spherical harmonics.

\begin{proposition}[Galerkin shape count]
  Fix $L$. After imposing algebraic slice conditions that remove the
  infinitesimal translation and rotation directions, suppose the resulting
  finite polynomial system has only isolated complex roots. Then the number of
  Galerkin shapes is bounded by the BKK mixed volume of that sliced system's
  Newton polytopes, and in particular by the Bezout bound $k^{N(L)}$ before
  using sparsity or symmetry.
\end{proposition}

\begin{remark}[Limit caveat]
  This is only a statement about the fixed-$L$ Galerkin model. It does not by
  itself imply finiteness of true shapes as $L\to\infty$: one still needs
  uniform a priori bounds, a no-escape compactness statement modulo $G$, and a
  stability result preventing spurious Galerkin roots or bifurcations in the
  high spherical-harmonic modes.
\end{remark}

\begin{remark}[How Galerkin can still help]
  A Galerkin computation becomes relevant to the transversality route if it is
  paired with uniform analytic estimates. The useful target is a high-mode
  bound of the form
  \[
    \|(I-P_L)T(\rho)\|_X \leq \varepsilon_L\,\Phi(\|\rho\|_X),
    \qquad \varepsilon_L\to 0,
  \]
  together with an inverse bound for the sliced finite-dimensional Jacobian at
  each computed root. A Newton--Kantorovich or Lyapunov--Schmidt argument could
  then validate an actual nearby solution and prove that no additional
  solutions occur in the corresponding slice neighborhood. This would support
  the compactness-plus-local-isolation proof, but only after the uniform
  high-mode estimate and the no-escape compactness statement are established.
\end{remark}

% -----------------------------------------------------------------------
\section{A Validated Direct Search Strategy for All Shapes}
% -----------------------------------------------------------------------

Assume that the transversality route above has been completed: modulo
$G=\mathbb R^3\rtimes SO(3)$ the non-trivial solution set in the chosen class
$\mathcal U\subset X$ is finite. The next goal is not just to know that the
set is finite, but to find every shape by a finite, checkable search. The
search should be designed so that failure of a required estimate is visible:
the algorithm must either certify all solutions, or report exactly which
analytic constant is missing.

\begin{remark}[Avoiding Galerkin truncation]
  The strategy in this section does not use a spherical-harmonic Galerkin
  cutoff. This avoids making the proof depend on a high-frequency tail constant
  $\tau_L$ for spherical harmonics. Instead we work directly on a mesh of
  $S^2$ and use a modulus of continuity, derivative bounds, and validated
  quadrature for the incidence-fiber integral. The finite-dimensional object is
  a rigorous enclosure of a compact class of functions, not an assertion that
  solutions lie in a finite-dimensional harmonic subspace.
\end{remark}

\subsection{Constants needed before the search}

The algorithm requires explicit versions of the estimates used in the
existence and finiteness proof.
\begin{enumerate}[leftmargin=2em]
  \item \textbf{A priori radius.} A bound
    \[
      \|\rho\|_X\leq R(c,k)
    \]
    for every non-trivial solution in the selected slice-normalised class.
    Ideally $R$ depends only on $c$ and $k$. In practice it may also depend on
    the chosen Sobolev index and on constants in the incidence-fiber estimate.

  \item \textbf{Non-collapse threshold.} A lower bound
    \[
      \|\rho\|_X\geq r(c,k)>0
    \]
    for every non-zero solution. A crude version may come from
    $\|\rho\|_X\leq |c|\,C\|\rho\|_X^k$, which gives
    $\|\rho\|_X\geq (|c|C)^{-1/(k-1)}$ whenever the multilinear estimate
    $\|T(\rho)\|_X\leq C\|\rho\|_X^k$ is valid.

  \item \textbf{Regularity and mesh modulus.} A derivative or modulus bound
    \[
      \|\rho\|_{C^m(S^2)}\leq B_m(c,k)
    \]
    for all solutions in the slice-normalised class. If the mesh size is $h$,
    this gives a computable interpolation envelope
    \[
      \|\rho-I_h\rho\|_{C^0}\leq \eta_h(B_m),
      \qquad \eta_h(B_m)\to0.
    \]
    This replaces the Galerkin tail constant. The mesh error is controlled by
    physical regularity on $S^2$, not by spherical-harmonic decay.

  \item \textbf{Validated quadrature for the fiber integral.} A computable
    error bound for evaluating
    \[
      T(\rho)(\omega)
      =
      \int_{V_k(\omega)}\rho(\omega_1)\cdots\rho(\omega_k)\,
      d\sigma_{\mathrm{fiber}}
    \]
    from interval enclosures of $\rho$ on the spherical mesh. Denote the
    resulting quadrature and geometry error by $q_h(B_m)$. It must tend to zero
    as the mesh is refined and must remain valid near singular or
    nearly-degenerate fibers.


  \item \textbf{Slice conditioning.} For each true solution $\rho$, a lower
    bound on the sliced linearisation
    \[
      \|L_\rho\eta\|_X\geq \mu(\rho)\|\eta\|_X,
      \qquad \eta\in S_\rho.
    \]
    To get a uniform count bound from packing, one wants
    $\mu_*:=\inf_\rho\mu(\rho)>0$. Compactness plus nondegeneracy gives
    $\mu_*>0$ abstractly; the search needs a computable lower bound.

  \item \textbf{Second derivative/Lipschitz bound.} A bound
    \[
      \|DF(\rho)-DF(\tilde\rho)\|_{S_\rho\to X}
      \leq M\|\rho-\tilde\rho\|_X
    \]
    on the radius-$R$ ball, after accounting for movement of the slice. This
    is needed for Newton--Kantorovich validation and for lower bounds on
    separation between distinct shapes.
\end{enumerate}

\subsection{From conditioning to separation of shapes}

The cleanest way to bound the number of shapes is a packing argument in the
quotient. Suppose we have a slice distance
\[
  d_G(\rho,\tilde\rho):=\inf_{g\in G}\|\rho-g\cdot\tilde\rho\|_X
\]
on the compact quotient of the solution set. If the sliced linearisation has
uniform inverse bound $\mu_*^{-1}$ and the derivative is $M$-Lipschitz, then
there is a quantitative isolation scale
\[
  \delta_* \simeq \frac{\mu_*}{M}
\]
such that two distinct shapes cannot lie in the same $\delta_*$-ball in the
quotient. More explicitly, if $\eta$ is a slice displacement from a solution
$\rho$, then
\[
  F(\rho+\eta)=L_\rho\eta+O(M\|\eta\|_X^2),
  \qquad
  \|F(\rho+\eta)\|_X
    \geq \mu_*\|\eta\|_X-CM\|\eta\|_X^2.
\]
Thus a second zero in the slice is impossible for
$0<\|\eta\|_X<\mu_*/(CM)$. This is the local separation estimate needed by the
search.

To turn this into a numerical upper bound, cover the slice-normalised compact
set
\[
  \mathcal K:=\{\rho:\ r(c,k)\leq\|\rho\|_X\leq R(c,k),\
      \|\rho\|_{C^m}\leq B_m,\ \rho\text{ satisfies the slice conditions}\}
\]
by mesh value boxes. If the triangulation has $N_h$ nodes, the naive effective
dimension is approximately the number of real node variables after imposing
the slice conditions. A very crude packing bound has the form
\[
  \#\{\text{shapes}\}
  \leq
  \left(1+\frac{2R_h}{\delta_*-2\eta_h-2q_h}\right)^{N_{\mathrm{mesh}}(h)},
\]
provided $\delta_*>2\eta_h+2q_h$. Here $R_h$ is the mesh-value radius,
$N_{\mathrm{mesh}}(h)$ is the number of real mesh degrees of freedom after
gauge conditions, and $q_h$ is the validated quadrature error. This is not
expected to be sharp. Its value is that it depends only on explicit constants
and certifies that the mesh is fine enough to separate distinct orbits.

\subsection{Validated fixed-point search on a spherical mesh}

The practical search should not enumerate all grid points blindly if this can
be avoided. A reasonable validated pipeline is:
\begin{enumerate}[leftmargin=2em]
  \item Choose a geodesic triangulation or quasi-uniform point set
    $\Omega_h=\{\omega_1,\ldots,\omega_{N_h}\}\subset S^2$ with mesh size $h$.
    Fix a gauge for translation and rotation by slice equations, for example
    orthogonality to the infinitesimal orbit directions.

  \item Represent a candidate function by intervals for its nodal values
    $(\rho(\omega_i))_{i=1}^{N_h}$ together with the global bounds
    $\|\rho\|_X\leq R$ and $\|\rho\|_{C^m}\leq B_m$. The interval box therefore
    represents all functions whose values lie in the box and whose oscillation
    between nodes is bounded by $\eta_h(B_m)$.

  \item For each box $B$, compute an interval enclosure for
    \[
      F(\rho)(\omega_i)=\rho(\omega_i)-cT(\rho)(\omega_i),
      \qquad i=1,\ldots,N_h,
    \]
    using validated quadrature over the incidence fibers. The enclosure must
    include both value uncertainty inside $B$ and the quadrature/mesh error
    $q_h(B_m)$.

  \item Discard $B$ if some node residual interval, enlarged by the modulus
    error needed to control points between nodes, excludes zero. This is the
    direct exclusion step.

  \item For a surviving box, apply an interval Newton or Krawczyk test to the
    finite nodal fixed-point map. The Jacobian used here is the discretized
    version of the sliced operator $L_\rho$, with interval bounds that account
    for mesh and quadrature errors.

  \item If the Krawczyk image lies strictly inside the box, validate a unique
    fixed point in that box. Then use the mesh residual estimate and the
    contraction/implicit-function estimate in the full space to lift the nodal
    fixed point to a unique true solution of $F(\rho)=0$ in the represented
    function class.

  \item Merge candidates whose slice distance is below the group-action
    tolerance; then explicitly minimise over rotations and translations to
    decide whether they represent the same shape.

  \item Certify completeness by showing that every grid box has either been
    excluded or assigned to the validation neighborhood of exactly one
    solution.
\end{enumerate}

This is a branch-and-bound search in the quotient, not a purely numerical
root-finder. The important feature is the exclusion certificate: without it,
finding many solutions gives no evidence that all shapes have been found.

\subsection{Main difficulties}

There are four serious obstructions.
\begin{enumerate}[leftmargin=2em]
  \item \textbf{Translation gauge.} The action
    $\rho(\omega)\mapsto e^{-ix\cdot\omega}\rho(\omega)$ is noncompact and
    can create rapid oscillation on $S^2$ as $|x|$ grows. Any finite search
    must either fix a canonical center/gauge or prove that large translations
    are impossible inside the chosen compact slice.

  \item \textbf{Uniform conditioning.} Nondegeneracy at each solution is not
    enough for an effective search. We need a computable lower bound
    $\mu_*>0$ for the sliced linearisation over all solutions. This may be the
    hardest constant to obtain without already knowing the solutions.

  \item \textbf{Mesh consistency.} The direct grid search is useful only if the
    modulus-of-continuity and quadrature errors are rigorous. The estimate must
    be strong enough that a nodal approximate zero cannot hide a large residual
    between mesh points or along a poorly resolved incidence fiber.

  \item \textbf{Dimension explosion.} Even after quotienting by $G$, the number
    of nodal variables grows like the number of mesh points on $S^2$. A raw
    grid search will be expensive; interval branch-and-bound, adaptive mesh
    refinement, fixed-point contraction tests, and symmetry decomposition
    should be used to reduce the search.
\end{enumerate}

The most realistic plan is therefore staged. First prove the analytic constants
\[
  R,\qquad r,\qquad B_m,\qquad \eta_h,\qquad q_h,\qquad \mu_*,\qquad M
\]
in a conservative norm. Then run interval fixed-point search on the sliced
spherical mesh. Finally validate each candidate, and certify that all remaining
boxes are empty.

% -----------------------------------------------------------------------
\section{\texorpdfstring{Restriction Constant Program on $S^1$}{Restriction Constant Program on S1}}
% -----------------------------------------------------------------------

The concrete target is the sharp Fourier extension constant on the unit circle
$S^1\subset\mathbb R^2$:
\[
  \mathcal R
  :=
  \sup_{0\neq f\in L^2(S^1)}
  \frac{\|\widehat{f\sigma}\|_{L^6(\mathbb R^2)}}{\|f\|_{L^2(S^1)}}.
\]
Here $\sigma$ is arclength measure on $S^1$ and
\[
  u(x)=\widehat{f\sigma}(x)
     =\int_{S^1} f(\omega)e^{ix\cdot\omega}\,d\sigma(\omega).
\]
Thus $\widehat u$ is supported on $S^1$ by construction. With the Laplacian
convention used in this note, every such $u$ satisfies
\[
  (-\Delta-1)u=0.
\]
If one writes $\Delta u=u$ instead, then one is using the opposite sign
convention for the Laplacian.

\subsection{Euler--Lagrange equation and normalization}

Normalize
\[
  \|f\|_{L^2(S^1)}=1.
\]
Then an extremizer satisfies the Euler--Lagrange equation
\[
  E^*(|Ef|^4Ef)=\Lambda f,
  \qquad
  \Lambda=\|Ef\|_{L^6(\mathbb R^2)}^6,
\]
where $E f=\widehat{f\sigma}$ and $E^*$ is the adjoint restriction operator.
Applying $E$ to both sides gives the nonlocal Helmholtz equation
\[
  (|u|^4u)*\widehat{\sigma}=\Lambda u,
  \qquad
  (-\Delta-1)u=0.
\]
Equivalently,
\[
  u=c\bigl((|u|^4u)*\widehat{\sigma}\bigr),
  \qquad c=\Lambda^{-1}.
\]
The absolute values are important: the restriction problem gives
$|u|^4u$, not generally $u^5$. The equation with $u^5$ is a useful model only
in a real nonnegative or otherwise reduced setting where $|u|^4u=u^5$.

The normalization is essential. Without it, if $u$ solves
$u=c((|u|^4u)*\widehat\sigma)$, then $\alpha u$ solves the same equation with
$c$ replaced by $c\alpha^{-4}$. Thus one cannot get a discrete list of
parameters $c$ until the amplitude is fixed, for example by
$\|f\|_2=1$.

\subsection{What ``finding all \texorpdfstring{$c$}{c}'' means}

After normalization, the search problem becomes the finite-dimensional-looking
but still analytic problem
\[
  \mathcal F(f,\Lambda)
  :=
  \bigl(E^*(|Ef|^4Ef)-\Lambda f,\ \|f\|_2^2-1\bigr)=0,
\]
modulo the symmetries of the circle: rotation of $S^1$, modulation/translation
of the extension, and constant phase. If the compactness and transversality
program proves that the normalized critical set modulo these symmetries is
finite, then the possible critical values
\[
  \Lambda_j=\|Ef_j\|_6^6
\]
are finite. The sharp restriction constant is then
\[
  \mathcal R=\max_j \Lambda_j^{1/6}.
\]
In the $c$-parametrization above, this is the same as finding the smallest
positive $c_j=\Lambda_j^{-1}$:
\[
  \mathcal R=c_{\min}^{-1/6}.
\]

\subsection{Direct certified search over critical values}

The direct search should therefore use $(f,\Lambda)$, or equivalently
$(f,c)$, as variables. A concrete certified program is:
\begin{enumerate}[leftmargin=2em]
  \item Prove a priori bounds for normalized critical points:
    \[
      \|f\|_{L^2}=1,\qquad
      \|f\|_{C^m(S^1)}\leq B_m,\qquad
      0<\Lambda_{\min}\leq\Lambda\leq\Lambda_{\max}.
    \]

  \item Fix gauges for the symmetries. For instance, impose a phase condition,
    a rotation condition on $S^1$, and a translation/modulation condition for
    $u=Ef$. These are the constraints that make two equivalent extremizers
    appear only once in the search.

  \item Mesh $S^1$ directly. Represent $f$ by interval boxes for its values on
    a circular grid, together with the derivative bound $B_m$, giving a
    rigorous interpolation error $\eta_h(B_m)$.

  \item On each box for $(f,\Lambda)$, enclose the residual
    \[
      E^*(|Ef|^4Ef)(\omega_i)-\Lambda f(\omega_i)
    \]
    at the mesh points, plus the normalization and gauge equations. The
    enclosure must include quadrature error for the $x$-space or convolution
    representation and interpolation error between grid points.

  \item Discard boxes whose interval residual excludes zero. For surviving
    boxes, apply an interval Newton or Krawczyk test to the augmented system.
    The linearization to bound is
    \[
      (h,\lambda)\mapsto
      E^*\!\left(4|u|^2\operatorname{Re}(\overline u\,Eh)\,u
        + |u|^4Eh\right)
      -\Lambda h-\lambda f,
    \]
    together with the derivative of the normalization and gauge constraints.

  \item Each validated box yields one critical point and one interval for
    $\Lambda_j$. Completeness is certified only after every box in the
    normalized search region is either discarded or assigned to exactly one
    validated critical point.
\end{enumerate}

At the end, sort the validated intervals for $\Lambda_j$. If the largest
interval is separated from the others, its sixth root gives the sharp
$L^2(S^1)\to L^6(\mathbb R^2)$ extension constant. If the top intervals overlap,
the mesh or conditioning estimates must be refined until the ordering is
certified.

\subsection{Main analytic bottlenecks for the circle problem}

The difficult estimates are now more specific than in the model problem:
\begin{enumerate}[leftmargin=2em]
  \item compactness of normalized critical points modulo the circle symmetries;
  \item a regularity bootstrap for critical points of
    $E^*(|Ef|^4Ef)=\Lambda f$;
  \item a quantitative slice theorem fixing phase, rotation, and translation;
  \item a lower bound on the augmented sliced linearization, uniform over all
    normalized critical points;
  \item validated quadrature or convolution estimates for the nonlinear
    residual on a direct mesh of $S^1$.
\end{enumerate}

The conclusion we want is not merely existence of an extremizer. We want a
finite certified list of normalized critical values $\Lambda_j$, because the
sharp constant is obtained by selecting the largest one.

% -----------------------------------------------------------------------
\section{Using Symmetry to Reduce to Finitely Many Critical Values}
% -----------------------------------------------------------------------

The role of symmetry is not, by itself, to prove finiteness. Rather, symmetry
removes the continuous families that are forced by invariance. The finiteness
principle should be formulated as
\[
  \begin{gathered}
  \text{compactness modulo symmetry}
  \quad+\quad
  \text{nondegeneracy on a slice}\\
  \Longrightarrow
  \text{finitely many critical values.}
  \end{gathered}
\]

\subsection{The symmetry group}

For the $S^1$ restriction problem the natural group is
\[
  G=U(1)\times SO(2)\ltimes\mathbb R^2.
\]
It acts on densities by
\[
  (e^{i\alpha},R,x_0)\cdot f(\omega)
  =
  e^{i\alpha}e^{-ix_0\cdot\omega}f(R^{-1}\omega).
\]
These transformations preserve $\|f\|_{L^2(S^1)}$ and
$\|Ef\|_{L^6(\mathbb R^2)}$, hence they preserve the Euler--Lagrange parameter
\[
  \Lambda=\|Ef\|_6^6.
\]
Thus every critical point generates a whole $G$-orbit with the same
$\Lambda$. The finite object is therefore not the raw critical set, but the
quotient
\[
  \mathcal Z
  :=
  \{(f,\Lambda):E^*(|Ef|^4Ef)=\Lambda f,\ \|f\|_2=1\}/G.
\]

\subsection{Compactness modulo symmetry}

The first target theorem is:
\[
  (f_n,\Lambda_n)\in\mathcal Z
  \quad\Longrightarrow\quad
  \exists\,g_n\in G
  \text{ such that } g_nf_n\to f_\infty
  \text{ strongly in the chosen space.}
\]
This is a concentration-compactness statement. It must rule out loss of mass,
splitting of mass, and escape through the translation/modulation parameter
$x_0$. The normalization $\|f_n\|_2=1$ rules out trivial scaling collapse, but
does not by itself rule out symmetry escape. That is exactly why the quotient
by $G$ is needed.

\subsection{Slice and infinitesimal symmetries}

At a critical point $f$, the infinitesimal symmetry directions are
\[
  if,\qquad
  \partial_\theta f,\qquad
  -i(a\cdot\omega)f,\quad a\in\mathbb R^2.
\]
The slice should impose orthogonality conditions against these directions, for
example in the $L^2(S^1)$ or $H^s(S^1)$ inner product. After adding the
normalization constraint, a nearby normalized critical point should be
representable uniquely, modulo any finite stabilizer, as
\[
  g\cdot(f+\eta),
  \qquad
  \eta\perp T_f(G\cdot f).
\]

\subsection{Nondegeneracy of the augmented equation}

Define the augmented Euler--Lagrange map
\[
  \mathcal F(f,\Lambda)
  :=
  \bigl(E^*(|Ef|^4Ef)-\Lambda f,\ \|f\|_2^2-1\bigr).
\]
The linearized operator in the variables $(f,\Lambda)$ is
\[
  \begin{aligned}
  D\mathcal F(f,\Lambda)[h,\lambda]
  &=
  \left(
    E^*\!\left(4|u|^2\operatorname{Re}(\overline u\,Eh)\,u
        + |u|^4Eh\right)
    -\Lambda h-\lambda f,\right.\\
  &\hspace{7em}\left.
    2\operatorname{Re}\langle f,h\rangle
  \right),
  \qquad u=Ef.
  \end{aligned}
\]
The required nondegeneracy condition is
\[
  \ker D\mathcal F(f,\Lambda)=T_f(G\cdot f).
\]
Equivalently, once one restricts to the slice and imposes the normalization,
the linearized operator has no kernel. This is the local discreteness
statement: the implicit function theorem then says that there are no nearby
critical points except those obtained by symmetry.

\subsection{Finiteness of admissible \texorpdfstring{$\Lambda_j$}{Lambda j}}

If the normalized critical set is compact modulo $G$, and if every normalized
critical point is nondegenerate modulo $G$, then $\mathcal Z$ is compact and
discrete. Therefore $\mathcal Z$ is finite. Since $\Lambda$ is constant along
each $G$-orbit, the set of admissible critical values
\[
  \{\Lambda_j\}
\]
is finite as well.

The sharp restriction constant is then obtained by selecting the largest
critical value:
\[
  \mathcal R=\max_j \Lambda_j^{1/6}.
\]
In the reciprocal parametrization
$u=c((|u|^4u)*\widehat\sigma)$, this is the same as selecting the smallest
positive $c_j=\Lambda_j^{-1}$.

\end{document}
